% ══════════════════════════════════════════════════════════════════════════════
%  LESSON 1 — WHAT IS SCIENCE?                                         (Week 1 · P1)
% ══════════════════════════════════════════════════════════════════════════════
\setcounter{section}{0}
\sectiondiv{What is Science?}{How we know what we know \textbullet{} Observations and inferences \textbullet{} Questions science can answer}\label{sec:1}
\lessoncodes{VC2S8H01 \textbullet{} VC2S8I01}

\vocabhead
\begin{ctab}
\begin{tabular}{|SW{3.4cm}|W{7.0cm}|}
  \hline\hd & \thead{Word} & \thead{Definition (copy from the slide)} \\ \hline
  & \scaf{\textbf{Science}} & \scaf{A way of finding out about the natural world by testing ideas against evidence.} \\ \hline
  & \textbf{Evidence} & \\ \hline
  & \textbf{Observation} & \\ \hline
  & \textbf{Inference} & \\ \hline
  & \textbf{Investigable question} & \\ \hline
\end{tabular}
\end{ctab}

% ─── PART A ───────────────────────────────────────────────────────────────────
\lessonpart{A}{Reading and comprehension}{Teal}\label{read:1}

\begin{readingbox}{Reading --- The Doctor Who Drank the Bacteria}
For most of the twentieth century, doctors were sure they knew what caused stomach ulcers: stress, spicy food and too much acid. Nothing, they said, could live in the acid of the stomach.

In Perth in the early 1980s, two Australian scientists looked more closely. Robin Warren, a pathologist, made a careful \textbf{observation}: when he looked at samples of ulcer patients' stomachs under a microscope, curved bacteria were there again and again. From this he made an \textbf{inference} --- perhaps the bacteria were causing the ulcers.

Most doctors did not believe it. So Warren and his colleague Barry Marshall set out to find \textbf{evidence}. They grew the bacteria in the lab and turned their idea into an \textbf{investigable question}: \emph{do these bacteria cause stomach illness?} Marshall even drank a sample. Within days he was ill, and the bacteria were found in his stomach.

Further tests showed that ulcers could be cured with antibiotics. In 2005 Marshall and Warren won the Nobel Prize. Their story shows how \textbf{science} works: ideas are tested against evidence, and even ideas everyone believes can change when new evidence comes in.
\end{readingbox}

\rcreset
\rcq{What did most doctors believe caused stomach ulcers?}{2.5cm}
\rcq{What observation did Robin Warren make?}{2.5cm}
\rcq{What inference did he make from it?}{2.5cm}
\rcq{Why did Marshall and Warren need to find evidence?}{2.5cm}
\rcq{Write the investigable question they tested.}{2.5cm}
\rcq{What does this story show about how scientific ideas can change?}{3cm}

% ─── PART B ───────────────────────────────────────────────────────────────────
\lessonpart{B}{Theory and examples}{Blue}\label{theory:1}

\subsection{Observations and Inferences}\label{sub:1.1}

\begin{skillbox}{Two different things}
An \textbf{observation} is something you notice with your senses or measure: \emph{the grass is wet}. An \textbf{inference} is an explanation of what you observed: \emph{it rained last night}. One observation can lead to several inferences --- the grass could also be wet from a sprinkler --- so scientists look for more evidence to decide between them.
\end{skillbox}

\begin{workedex}{Observation or inference?}
A student finds a glass on the bench with a little water in the bottom. Sort the statements.

\smallskip
\emph{There is water in the glass.} \textbf{Observation} --- it can be seen.\\
\emph{The glass is cold to touch.} \textbf{Observation} --- it can be felt.\\
\emph{Someone drank most of the water.} \textbf{Inference} --- it explains what is seen.
\end{workedex}

\begin{wedo}{Observation or inference?}
You see muddy footprints leading from the back door to the fridge.
\pt{a}Write one observation.
\answerbox{2.5cm}
\pt{b}Write one inference.
\answerbox{2.5cm}
\end{wedo}

\subsection{Questions Science Can Answer}\label{sub:1.2}

\begin{skillbox}{An investigable question}
Science can only answer a question that can be tested by collecting evidence. \emph{Which ice-cream flavour is best?} is a matter of opinion. An \textbf{investigable question} names what you will change and what you will measure:
\par\smallskip\centerline{\textbf{How does} [the thing I change] \textbf{affect} [the thing I measure]\textbf{?}}
\end{skillbox}

\begin{workedex}{Making a question investigable}
\emph{Do plants like music?} is hard to test --- what does ``like'' mean?

\smallskip
\textbf{Investigable:} \emph{How does playing music affect the height a bean plant grows in two weeks?} It names what is changed (music) and what is measured (height).
\end{workedex}

\begin{wedo}{Making a question investigable}
Rewrite \emph{Is sugar good at dissolving?} as an investigable question.
\answerbox{3cm}
\end{wedo}

% ─── PART C ───────────────────────────────────────────────────────────────────
\lessonpart{C}{Practice}{Amber}\label{prac:1}

\task Write \emph{O} for observation or \emph{I} for inference beside each statement.
\begin{ctab}
\begin{tabular}{|SW{8.4cm}|Z{1.6cm}|}
  \hline\hd & \thead{Statement} & \thead{O or I?} \\ \hline
  & \scaf{The candle flame is yellow.} & \scaf{O} \\ \hline
  & The candle is getting shorter. & \\ \hline
  & The wax is burning as fuel. & \\ \hline
  & The air above the flame feels hot. & \\ \hline
  & The flame needs air to keep burning. & \\ \hline
\end{tabular}
\end{ctab}

\needspace{12\baselineskip}
\task A lunchbox left in a locker for two weeks is opened. The bread has green, fuzzy spots, and it smells sour.
\pt{a}Write two observations.
\answerbox{3.5cm}
\pt{b}Write one inference, and say what further evidence would support it.
\answerbox{5cm}

\needspace{12\baselineskip}
\task Tick the questions that science could investigate.
\begin{ctab}
\begin{tabular}{|SW{8.4cm}|Z{1.6cm}|}
  \hline\hd & \thead{Question} & \thead{Testable?} \\ \hline
  & \scaf{Does salt make ice melt faster?} & \scaf{\ding{51}} \\ \hline
  & Which colour is the prettiest? & \\ \hline
  & How does the length of a pendulum affect its swing time? & \\ \hline
  & Is football better than netball? & \\ \hline
  & Does warm water dissolve sugar faster than cold water? & \\ \hline
\end{tabular}
\end{ctab}

\needspace{10\baselineskip}
\task Rewrite each question as an investigable question.
\pt{a}Do paper planes fly well?
\answerbox{3cm}
\pt{b}Is exercise good for your heart?
\answerbox{3cm}

\needspace{10\baselineskip}
\task Your teacher shows you a mystery object. Write one observation and one inference about it.
\answerbox{5cm}

% ─── PART D ───────────────────────────────────────────────────────────────────
\lessonpart{D}{Above level}{Violet}\label{ext:1}

\begin{yr8box}
These questions go beyond the Year 7 standard.

\medskip
\begin{keepext}
\exttask Most doctors did not accept Marshall and Warren's idea at first. Explain why scientists are right to ask for strong evidence before they change an accepted idea, and why they must still be willing to change it.
\answerbox{5cm}
\end{keepext}
\begin{keepext}
\exttask A website says that wearing a copper bracelet cures sore joints. Describe how a scientist could test this claim fairly.
\answerbox{5cm}
\end{keepext}
\end{yr8box}

% ══════════════════════════════════════════════════════════════════════════════
%  LESSON 4 — THE SCIENTIFIC METHOD                                     (Week 2 · P1)
% ══════════════════════════════════════════════════════════════════════════════
\setcounter{section}{3}
\sectiondiv{The Scientific Method}{The steps of an investigation \textbullet{} Independent, dependent and controlled variables \textbullet{} Writing a hypothesis}\label{sec:4}
\lessoncodes{VC2S8I01 \textbullet{} VC2S8I02 \textbullet{} VC2S8H01}

\vocabhead
\begin{ctab}
\begin{tabular}{|SW{3.4cm}|W{7.0cm}|}
  \hline\hd & \thead{Word} & \thead{Definition (copy from the slide)} \\ \hline
  & \scaf{\textbf{Hypothesis}} & \scaf{A prediction that can be tested, with a reason: If \dots\ then \dots\ because \dots} \\ \hline
  & \textbf{Independent variable} & \\ \hline
  & \textbf{Dependent variable} & \\ \hline
  & \textbf{Controlled variable} & \\ \hline
  & \textbf{Fair test} & \\ \hline
  & \textbf{Conclusion} & \\ \hline
\end{tabular}
\end{ctab}

% ─── PART A ───────────────────────────────────────────────────────────────────
\lessonpart{A}{Reading and comprehension}{Teal}\label{read:4}

\begin{readingbox}{Reading --- The Paper Helicopter Challenge}
Ms Tran's class was given a challenge: make a paper helicopter that takes as long as possible to fall. Mia noticed that helicopters with longer blades seemed to spin down more slowly. She asked a question: \emph{how does the length of the blades affect the time a paper helicopter takes to fall?}

Mia wrote a \textbf{hypothesis}: \emph{If the blades are longer, then the helicopter will take longer to fall, because longer blades push against more air.}

To test it she changed only the blade length: 4, 6, 8 and 10\,cm. This was her \textbf{independent variable}. She measured the fall time with a stopwatch --- her \textbf{dependent variable}. Everything else stayed the same: the same paper, the same paper clip on the bottom, the same drop height of 2 metres. These were her \textbf{controlled variables}. Because only one thing changed, her experiment was a \textbf{fair test}.

She timed each helicopter three times and averaged the results. The 10\,cm helicopter took longest to fall. Her \textbf{conclusion}: longer blades made the helicopter fall more slowly, so her hypothesis was supported.
\end{readingbox}

\rcreset
\rcq{What question did Mia investigate?}{2.5cm}
\rcq{What was her hypothesis?}{2.5cm}
\rcq{Name her independent variable and her dependent variable.}{2.5cm}
\rcq{Name two of her controlled variables.}{2.5cm}
\rcq{Why was her experiment a fair test?}{2.5cm}
\rcq{Why did she time each helicopter three times?}{2.5cm}

% ─── PART B ───────────────────────────────────────────────────────────────────
\lessonpart{B}{Theory and examples}{Blue}\label{theory:4}

\subsection{The Steps of an Investigation}\label{sub:4.1}

\begin{center}
\begin{tikzpicture}[font=\small, box/.style={draw=Ink, line width=0.8pt, rounded corners=1.5mm, minimum width=19mm, minimum height=9mm, align=center, fill=white}]
  \node[box] (q) at (0,0) {Question};
  \node[box] (h) at (2.35,0) {Hypothesis};
  \node[box] (m) at (4.7,0) {Method};
  \node[box] (r) at (7.05,0) {Results};
  \node[box] (c) at (9.4,0) {Conclusion};
  \foreach \a/\b in {q/h,h/m,m/r,r/c}{\draw[Ink, line width=1pt, -{Stealth[length=5pt]}] (\a) -- (\b);}
  \draw[Ink, line width=0.8pt, dashed, -{Stealth[length=5pt]}] (c.south) -- ++(0,-0.6) -| (q.south) node[pos=0.25, below, font=\scriptsize] {new questions};
\end{tikzpicture}
\end{center}

\begin{skillbox}{What each step does}
\textbf{Question:} what you want to find out. \textbf{Hypothesis:} your prediction and the reason for it. \textbf{Method:} numbered steps someone else could repeat. \textbf{Results:} measurements in a table, and often a graph. \textbf{Conclusion:} what the results show, and whether they support the hypothesis.
\end{skillbox}

\subsection{Variables}\label{sub:4.2}

\begin{skillbox}{Three kinds of variable}
The \textbf{independent variable} is what you change. The \textbf{dependent variable} is what you measure. The \textbf{controlled variables} are what you keep the same, so the test is fair.
\end{skillbox}

\begin{workedex}{Naming variables}
A student tests how the temperature of water affects how fast a sugar cube dissolves.

\smallskip
\textbf{Independent:} the temperature of the water.\\
\textbf{Dependent:} the time for the cube to dissolve.\\
\textbf{Controlled:} the volume of water, the size of the sugar cube, stirring (or not).
\end{workedex}

\begin{wedo}{Naming variables}
A student tests how the height of a ramp affects how far a toy car rolls.
\pt{a}Independent variable:
\answerbox{1.2cm}
\pt{b}Dependent variable:
\answerbox{1.2cm}
\pt{c}Two controlled variables:
\answerbox{2.5cm}
\end{wedo}

\subsection{Writing a Hypothesis}\label{sub:4.3}

\begin{workedex}{A hypothesis}
\textbf{Question:} How does the temperature of water affect how fast a sugar cube dissolves?

\smallskip
\textbf{Hypothesis:} \emph{If the water is hotter, then the sugar cube will dissolve faster, because hot water has more energy to break the sugar apart quickly.}
\end{workedex}

\begin{wedo}{A hypothesis}
Write a hypothesis for the toy-car ramp investigation.
\answerbox{4cm}
\end{wedo}

% ─── PART C ───────────────────────────────────────────────────────────────────
\lessonpart{C}{Practice}{Amber}\label{prac:4}

\task Number the steps of an investigation in order, 1 to 5.
\begin{ctab}
\begin{tabular}{|SW{8.4cm}|Z{1.6cm}|}
  \hline\hd & \thead{Step} & \thead{Order} \\ \hline
  & Results --- record measurements in a table & \\ \hline
  & Question --- decide what to find out & \\ \hline
  & Conclusion --- say what the results show & \\ \hline
  & Method --- write the steps to follow & \\ \hline
  & Hypothesis --- predict, with a reason & \\ \hline
\end{tabular}
\end{ctab}

\needspace{14\baselineskip}
\task Identify the variables in each investigation.
\begin{ctab}
\begin{tabular}{|SW{3.6cm}|W{2.2cm}|W{2.2cm}|W{2.4cm}|}
  \hline\hd & \thead{Investigation} & \thead{Independent} & \thead{Dependent} & \thead{One controlled} \\ \hline
  & \scaf{Does the size of a parachute change its fall time?} & \scaf{Size of parachute} & \scaf{Fall time} & \scaf{Drop height} \\ \hline
  & How does the amount of light affect plant growth? & & & \\ \hline
  & How does the mass of a ball affect how high it bounces? & & & \\ \hline
  & Does the type of surface change how far a marble rolls? & & & \\ \hline
\end{tabular}
\end{ctab}

\needspace{10\baselineskip}
\task Write a hypothesis for this question: \emph{How does the length of a pendulum affect the time of one swing?}
\answerbox{5cm}

\needspace{14\baselineskip}
\task A class drops paper helicopters to see how the number of paper clips on the bottom affects the fall time.
\pt{a}Name the independent and dependent variables.
\answerbox{3cm}
\pt{b}Name two variables they must control, and explain why.
\answerbox{5cm}

% ─── PART D ───────────────────────────────────────────────────────────────────
\lessonpart{D}{Above level}{Violet}\label{ext:4}

\begin{yr8box}
These questions go beyond the Year 7 standard.

\medskip
\begin{keepext}
\exttask Sam tests two paper helicopters. One has longer blades \emph{and} is made from thicker paper. It falls more slowly. Explain why Sam cannot conclude that longer blades make a helicopter fall more slowly, and how to fix the test.
\answerbox{5cm}
\end{keepext}
\begin{keepext}
\exttask Mia's results did not support a classmate's hypothesis. Explain why a hypothesis that turns out wrong is still useful in science.
\answerbox{5cm}
\end{keepext}
\end{yr8box}

% ══════════════════════════════════════════════════════════════════════════════
%  LESSON 5 — SAFETY                                                    (Week 2 · P2)
% ══════════════════════════════════════════════════════════════════════════════
\setcounter{section}{4}
\sectiondiv{Safety}{Lab rules and why they matter \textbullet{} Hazards, risks and controls \textbullet{} Hazard symbols \textbullet{} A simple risk assessment}\label{sec:5}
\lessoncodes{VC2S8I02}

\vocabhead
\begin{ctab}
\begin{tabular}{|SW{3.4cm}|W{7.0cm}|}
  \hline\hd & \thead{Word} & \thead{Definition (copy from the slide)} \\ \hline
  & \scaf{\textbf{Hazard}} & \scaf{Anything that could cause harm, such as a flame, hot water or broken glass.} \\ \hline
  & \textbf{Risk} & \\ \hline
  & \textbf{Control} & \\ \hline
  & \textbf{Risk assessment} & \\ \hline
  & \textbf{Hazard symbol} & \\ \hline
  & \textbf{PPE} & \\ \hline
\end{tabular}
\end{ctab}

% ─── PART A ───────────────────────────────────────────────────────────────────
\lessonpart{A}{Reading and comprehension}{Teal}\label{read:5}

\begin{readingbox}{Reading --- Why Labs Have Rules}
A school laboratory holds things you will not find in most classrooms: gas taps, open flames, glassware, hot plates and chemicals. Each of these is a \textbf{hazard}: something that could cause harm. A hazard is not the same as a \textbf{risk}. The risk is the chance that the hazard actually hurts someone, and how badly. A bottle of acid locked in a cupboard is a hazard with almost no risk; the same acid in an open beaker at the edge of a crowded bench is a high risk.

Scientists reduce risk with \textbf{controls}: actions or equipment that make harm less likely. Wearing safety glasses, tying back long hair and standing up while heating are controls. Safety glasses, gloves and lab coats are called personal protective equipment, or \textbf{PPE}.

Before an experiment, scientists write a \textbf{risk assessment}. For each hazard they list the risk and the control. Chemical containers carry \textbf{hazard symbols}: red-bordered diamonds that warn, for example, that a substance is flammable or an irritant.

Most lab accidents happen when someone rushes, ignores an instruction or does not report a small problem. Following the rules every time keeps you and everyone around you safe.
\end{readingbox}

\rcreset
\rcq{Give three hazards found in a school laboratory.}{2.5cm}
\rcq{What is the difference between a hazard and a risk?}{2.5cm}
\rcq{Explain why the acid on the edge of a crowded bench is a higher risk than the acid in the cupboard.}{2.5cm}
\rcq{What is a control? Give two examples.}{2.5cm}
\rcq{What does PPE stand for? Give an example.}{2.5cm}
\rcq{How do most lab accidents happen?}{2.5cm}

% ─── PART B ───────────────────────────────────────────────────────────────────
\lessonpart{B}{Theory and examples}{Blue}\label{theory:5}

\subsection{Hazard, Risk and Control}\label{sub:5.1}

\begin{workedex}{Hazard, risk and control}
A student heats water in a beaker over a Bunsen burner.

\smallskip
\textbf{Hazard:} hot water and hot glassware.\\
\textbf{Risk:} burns or scalds if the beaker is touched or knocked over.\\
\textbf{Control:} stand up while heating, keep the beaker in the middle of the gauze, and let it cool before touching it.
\end{workedex}

\begin{wedo}{Hazard, risk and control}
A student cuts a potato into cubes with a scalpel.
\pt{a}Hazard:
\answerbox{1.2cm}
\pt{b}Risk:
\answerbox{2.5cm}
\pt{c}Control:
\answerbox{2.5cm}
\end{wedo}

\subsection{Hazard Symbols}\label{sub:5.2}

\begin{center}
\newcommand{\haz}[2]{\begin{tikzpicture}[baseline=0]
  \draw[Ink, line width=1.6pt, fill=white] (0,0.95) -- (0.95,0) -- (0,-0.95) -- (-0.95,0) -- cycle;
  #1
  \node[font=\scriptsize\bfseries, anchor=north, align=center] at (0,-1.05) {#2};
\end{tikzpicture}}
\haz{\fill[Ink] (0,-0.5) .. controls (-0.45,-0.35) and (-0.3,0.1) .. (-0.05,0.55) .. controls (0.05,0.2) and (0.15,0.1) .. (0.2,0.35) .. controls (0.45,0) and (0.4,-0.35) .. (0,-0.5);
     \draw[Ink, line width=0.8pt] (-0.45,-0.6) -- (0.45,-0.6);}{Flammable}\hspace{6mm}%
\haz{\fill[Ink] (-0.08,-0.1) rectangle (0.08,0.5); \fill[Ink] (0,-0.35) circle (0.09);}{Irritant\\or harmful}\hspace{6mm}%
\haz{\fill[Ink] (0,0) circle (0.14);
     \foreach \a in {0,45,...,315}{\draw[Ink, line width=1.2pt] (\a:0.22) -- (\a:0.5);}}{Explosive}\hspace{6mm}%
\haz{\draw[Ink, line width=1pt] (-0.4,0.35) -- (-0.1,0.15) (0.4,0.35) -- (0.1,0.15);
     \fill[Ink] (-0.1,0.1) circle (0.05) (0.1,0.1) circle (0.05);
     \fill[Ink] (-0.5,-0.55) rectangle (-0.05,-0.42) (0.05,-0.55) rectangle (0.5,-0.42);
     \draw[Ink, line width=0.8pt] (-0.15,0.0) -- (-0.25,-0.35) (0.15,0.0) -- (0.25,-0.35);}{Corrosive}
\end{center}
{\small\itshape On real labels these diamonds have a red border. Corrosive substances eat away skin and metal.}

\begin{workedex}{Reading a label}
A bottle of methylated spirits shows the \emph{flammable} symbol. \textbf{What it means:} it catches fire easily. \textbf{Control:} keep it away from flames, and put the lid back on straight away.
\end{workedex}

\begin{wedo}{Reading a label}
A bottle of dilute hydrochloric acid shows the \emph{corrosive} symbol. Explain what this means and give one control.
\answerbox{3.5cm}
\end{wedo}

\subsection{A Simple Risk Assessment}\label{sub:5.3}

\begin{skillbox}{Writing a risk assessment}
List every hazard in the experiment. For each one, write the risk (what could happen) and the control (how you will stop it). A control is something you \emph{do} --- ``be careful'' is not a control.
\end{skillbox}

% ─── PART C ───────────────────────────────────────────────────────────────────
\lessonpart{C}{Practice}{Amber}\label{prac:5}

\task Spot the hazards. In this lab, a student has a bag on the floor in the aisle, long hair hanging loose, and safety glasses pushed up on their head. A beaker of water sits right at the edge of the bench, a Bunsen burner has been left on with a yellow flame while nobody uses it, and a broken test tube is in the paper bin. Complete the table for four of the hazards.
\begin{ctab}
\begin{tabular}{|SW{4.0cm}|W{6.2cm}|}
  \hline\hd & \thead{Hazard} & \thead{What should be done} \\ \hline
  & \scaf{Bag on the floor in the aisle} & \scaf{Put it under the bench so no one trips.} \\ \hline
  & & \\ \hline
  & & \\ \hline
  & & \\ \hline
  & & \\ \hline
\end{tabular}
\end{ctab}

\needspace{14\baselineskip}
\task Complete a risk assessment for this experiment: \emph{heating 100\,mL of salt water in an evaporating basin over a Bunsen burner until crystals form.}
\begin{ctab}
\begin{tabular}{|SW{2.8cm}|W{3.4cm}|W{3.8cm}|}
  \hline\hd & \thead{Hazard} & \thead{Risk} & \thead{Control} \\ \hline
  & \scaf{Bunsen flame} & \scaf{Burns; hair or paper catching fire} & \scaf{Tie hair back; use the safety flame when not heating.} \\ \hline
  & Hot evaporating basin & & \\ \hline
  & Spitting salt solution & & \\ \hline
  & Gas & & \\ \hline
\end{tabular}
\end{ctab}

\needspace{10\baselineskip}
\task Explain why ``be careful'' is not a control. Give a better control for carrying a beaker of hot water.
\answerbox{5cm}

\needspace{10\baselineskip}
\task Choose three of the lab rules in your safety contract. For each, explain the hazard it protects you from.
\answerbox{5cm}

% ─── PART D ───────────────────────────────────────────────────────────────────
\lessonpart{D}{Above level}{Violet}\label{ext:5}

\begin{yr8box}
These questions go beyond the Year 7 standard.

\medskip
\begin{keepext}
\exttask A teacher says: ``We cannot remove every hazard from a lab, but we can make the risk very small.'' Explain what she means, using heating water as an example.
\answerbox{5cm}
\end{keepext}
\begin{keepext}
\exttask Some chemicals are only used by teachers, in a fume cupboard. Suggest why, and what kind of hazard a fume cupboard controls.
\answerbox{5cm}
\end{keepext}
\end{yr8box}
